\documentclass[11pt,a4paper]{article}
\usepackage[T1]{fontenc}\usepackage{lmodern,amsmath,amssymb,amsthm,mathtools,microtype}
\usepackage[margin=25mm]{geometry}\usepackage[colorlinks=true,linkcolor=blue,urlcolor=blue,citecolor=blue]{hyperref}
\newtheorem{theorem}{Theorem}[section]\newtheorem{lemma}[theorem]{Lemma}\newtheorem{corollary}[theorem]{Corollary}
\newcommand{\Q}{\mathbb Q}\newcommand{\Z}{\mathbb Z}\DeclareMathOperator{\Tr}{Tr}
\setlength{\emergencystretch}{2em}
\title{Heights of positive weighted elliptic traces}
\author{}\date{14 November 2026}
\begin{document}\maketitle
\begin{abstract}
A positive weighted trace of several elliptic abscissas controls all their
indices, even when the weights vary through rational squares. The bound
is linear in the heights of the weights and the resulting trace. The
additive constant is ineffective. The result includes half-integral
indices, repeated points, negative indices, and finite torsion points.
The proof checks the exceptional translates and projective embedding in
the published family height theorem. A separate quadratic-reciprocity
lemma shows that the existing Gaussian prime conditions make the chosen
half-generator divisible by two in both residue groups. Together these
results permit two point differences to replace two stored doubled
points in the quartic construction.
\end{abstract}

\section{The curve, indices, and statement}
Put $K=\Q(\rho)$, $\rho^2=2$, $d=3-\rho$, $d'=3+\rho$,
$D=d^2$, and
\[
 E:y^2=x^3-Dx,\qquad R_0=(1+\rho,2\rho),\qquad P=2R_0.
\]
Retain the arithmetic hypothesis $2E(K)=\Z P$, with $P$ of infinite
order. Every $Q\in E(K)$ has a unique expression
$Q=[m]R_0+T_0$, where $m\in\Z$ and $T_0\in E[2](K)$.
Indeed, $2Q=mP$ implies $2(Q-mR_0)=O$, and uniqueness follows from
non-torsion. Define its index by $\nu(Q)=m/2$.
Write $h$ for absolute logarithmic Weil height, and $\widehat h$ for
elliptic canonical height. Thus
$\widehat h(Q)=\nu(Q)^2\widehat h(P)$.

\begin{theorem}\label{thm:weighted}
Fix a positive integer $s$ and fixed polynomials
$w_1,\ldots,w_s\in\Q[Z]$. There are constants $A>0,J\ge0$ such
that, for every $r\in\Q$ with $w_i(r)>0$ for all $i$, and every
$Q_i\in E(K)\setminus\{O\}$, the number
\[
 S=\frac14\Tr_{K/\Q}\sum_{i=1}^s w_i(r)x(Q_i)
\]
satisfies
\[
 \max_i|\nu(Q_i)|^2\le A\bigl(h(r)+h(S)\bigr)+J.
\]
The constants depend only on the fixed curve, subgroup and polynomials,
not on $r$ or the points. The argument does not make $J$ effective.
In particular the conclusion applies to any prescribed assignment of
weights $1$ and $r^2$, on $r\ne0$.
\end{theorem}

This is a theorem relative to the displayed group hypothesis, not a new
$abc$ assumption. The analytic input to its proof is the published
Theorem~1.1 of von Buhren~\cite{VB}. No prime-pattern or pole-realization
hypothesis enters this height argument. The positivity is imposed at the
actual rational parameter, not at every algebraic point of its projective
parameter space.

\section{Independence of the conjugate generators}
Let $E_0:Y^2=X^3-X$. Over a fixed field containing $K(i,\sqrt d,\sqrt{d'})$,
use the pointed isomorphisms
\[
 (x,y)\longmapsto (x/d,y/d^{3/2}),\qquad
 (x,y)\longmapsto (x/d',y/(d')^{3/2}).
\]
Let $R,R'$ be the respective images of $R_0$ and its conjugate.

\begin{lemma}\label{lem:independent}
The points $R,R'$ are linearly independent modulo torsion over
$\operatorname{End}(E_0)=\Z[i]$.
\end{lemma}
\begin{proof}
The square classes of $d,d'$ are independent in $K(i)^\times/K(i)^{\times2}$.
For completeness, an element of $K$ is a square in $K(i)$ only if it or
its negative is a square in $K$: the imaginary part of $(a+bi)^2$
forces $ab=0$. Neither $d$ nor $-d$ is a square in $K$, because its
norm is 7, not a rational square. The same holds for $d'$. Their product
is 7. If $(a+b\rho)^2=7$ with $a,b\in\Q$, then $ab=0$, giving either
$a^2=7$ or $2b^2=7$, both impossible. The negative of 7 is not a square
in the real field $K$. This proves the claimed independence.

There is consequently an automorphism of the fixed compositum fixing
$K(i,\sqrt d)$ and sending $\sqrt{d'}$ to $-\sqrt{d'}$. It fixes $R$
and negates $R'$, while fixing every endomorphism in $\Z[i]$. If
$\alpha R+\beta R'$ is torsion, subtraction of its image gives
$2\beta R'$ torsion. A nonzero elliptic endomorphism is an isogeny;
its inverse image of torsion consists of torsion points. Since $R'$ is
non-torsion, $\beta=0$, and then $\alpha=0$.
\end{proof}

\section{Noncancellation at poles}
Consider an actual tuple of points of $E_0^{2k}$ whose first block has
coordinates $[m_j]R+T_j$ and whose second has coordinates
$[m_j]R'+T_j'$, where $m_j\ne0$ and the $T_j,T_j'$ are torsion. Fix
positive rational weights $a_j$. The function to be considered is
\[
 F=\frac d4\sum_{j=1}^k a_jx_j+
   \frac {d'}4\sum_{j=1}^k a_jx_j'+c,
\]
where $c$ is any fixed constant.

\begin{lemma}\label{lem:translates}
No positive-dimensional abelian translate through this tuple, meeting
the domain of $F$, has constant restriction of $F$.
\end{lemma}
\begin{proof}
Any positive-dimensional abelian subvariety of a power of $E_0$ contains
an elliptic subvariety isogenous to $E_0$. This follows from complete
reducibility and the fact that every simple constituent of a subvariety
of $E_0^{2k}$ is isogenous to $E_0$. Compose an isogeny from $E_0$ to
such a subvariety with the inclusion. The resulting nonzero direction
has coordinate endomorphisms $\alpha_j,\alpha_j'\in\Z[i]$.
If a translate of the larger subvariety had constant $F$, the translate
in this direction through our tuple would also have constant $F$.

At a pole of a nonconstant coordinate in the first block, two such
coordinates $i,j$ can have a common pole only if
$[\alpha_j]Q_i-[\alpha_i]Q_j=O$. Non-torsion then gives
$m_i\alpha_j=m_j\alpha_i$. In a local analytic parameter on $E_0$,
the principal part of $x$ at the identity is $z^{-2}$. Thus the
coefficient of the common double pole, summed over all first-block
coordinates having that pole, is
\[
 \frac d{4\alpha_i^2}\sum_j a_j(m_i/m_j)^2\ne0.
\]
The sum is a strictly positive rational number. This argument allows
negative or repeated indices and arbitrary torsion translates.

A first-block and a second-block nonconstant coordinate cannot have
a common pole: the same elimination would give a relation
$m_i\alpha_j'R-m_j\alpha_iR'$ modulo torsion, contradicting
Lemma~\ref{lem:independent}. The identical positive-sum argument holds
within the second block. At least one direction coordinate is nonzero,
so at least one nonconstant summand has a pole, and it cannot cancel.
Consequently the restriction is not constant.
\end{proof}

\section{Application of the family height theorem}
\begin{proof}[Proof of Theorem~\ref{thm:weighted}]
First separate the finitely many possibilities for those $Q_i$ whose
index is zero. Such points are finite two-torsion points; freeze their
coordinates and omit them from the abelian ambient space. Their
contributions to the function are fixed polynomials in $r$. There are
only finitely many choices of omitted index sets and torsion points.
If all points are omitted, the left side of the asserted bound is zero.
Otherwise let $k\ge1$ denote the number of remaining indices.

On $E_0^{2k}$, with an additional parameter $r$, form the rational
function
\[
 F_r=\frac d4\sum_j w_j(r)x_j+
       \frac {d'}4\sum_j w_j(r)x_j'+c(r),
\]
where $c(r)$ accounts for the frozen torsion coordinates. Let
$\mathcal V$ be the integral reduced closure of its graph in
$E_0^{2k}\times(\mathbf P^1_r\times\mathbf P^1_S)$. This closure is
integral because the graph over its nonempty rational-function domain
is integral. All actual tuples under consideration lie in that domain:
$r$ and all abscissas are finite.

Fix the finite-rank group generated by $R$ in the first-block factors,
$R'$ in the second-block factors, and the finitely many needed torsion
coordinates. It is independent of $r$ and of the integer indices.
At an actual parameter with all weights positive,
Lemma~\ref{lem:translates} excludes every positive-dimensional abelian
translate in the fiber through our tuple. A translate in the projective
graph fiber through a domain point would meet the domain densely and
make the displayed rational function constant there; thus graph
boundary components do not evade the lemma.

We check all additional hypotheses of \cite[Theorem~1.1, p.~308]{VB}.
Embed each $E_0$ by $\mathcal O(3O)$ and its product by Segre. For
$n\ge1$, the plane-cubic sequence
\[
0\to\mathcal O_{\mathbf P^2}(n-3)\to\mathcal O_{\mathbf P^2}(n)
 \to\mathcal O_{E_0}(3nO)\to0
\]
has surjective restriction of global sections, since
$H^1(\mathbf P^2,\mathcal O(n-3))=0$. Tensoring these surjections
and using Segre gives the source's condition (1.1) for every $n\ge1$.
The resulting line bundle on the abelian variety is ample and symmetric.
Embed the parameter surface by $\mathcal O(1,1)$, whose height is
$h(r)+h(S)$ up to a fixed bounded term. The ambient parameter variety
is projective and integral. Every point of the chosen finite-rank group
belongs to the source set $C(\Gamma,\epsilon)$ by taking the error point
to be zero. The source theorem therefore bounds its canonical product
height by a fixed multiple of $h(r)+h(S)$ plus a fixed additive constant.

That canonical height is a fixed positive multiple of
$2\sum_j\widehat h(Q_j)$, by pointed-isomorphism and Galois invariance.
As $\widehat h(P)>0$, it bounds the largest squared index. Take the
largest constants over the finitely many frozen-torsion cases. All
constants are now independent of $r$ and the points. The cited theorem
does not make the additive constant effective.
\end{proof}

\begin{corollary}\label{cor:four}
For finite $Q_1,Q_3\in E(K)$ and nonzero doubled points
$Q_2=[a]P,Q_4=[b]P$, with $a,b\in\Z\setminus\{0\}$, put
$\nu(Q_1)=n$, $2\nu(Q_3)=t\in\Z$ and
\[
 T=\frac14\Tr\bigl(x(Q_1)+x(Q_3)+r^2(x(Q_2)+x(Q_4))\bigr).
\]
There are fixed $A,J$ with
\[
 \max(|n|,|a|,|t|,|b|)^2\le A(h(r)+h(T))+J\qquad(r\ne0).
\]
If $r>1$, then $T>0$.
\end{corollary}
\begin{proof}
The factor two in $t$ only enlarges the height constant by four.
At either real embedding every finite elliptic point has $x\ge-d$,
since $x(x-d)(x+d)\ge0$. A nonzero doubled point has square abscissa
and is non-torsion, hence $x>d$. Summing conjugate embeddings gives
$\Tr(x(Q_1)+x(Q_3))\ge-12$ and
$\Tr(x(Q_2)+x(Q_4))>12$, proving $4T>12r^2-12$.
\end{proof}

\section{Divisibility of the half-generator at selected primes}
\begin{lemma}\label{lem:gaussian}
Let $p=A^2+B^2$ be prime, with $A\equiv9\pmod {16}$ and
$B\equiv0\pmod {16}$. Then $R_0$ and its conjugate belong to
$2E(\mathbf F_p)$ in their respective reductions. Consequently
$\#E(\mathbf F_p)/4$ annihilates the reduction of $R_0$ whenever
$E$ has good reduction, and the analogous statement holds for the
conjugate curve.
\end{lemma}
\begin{proof}
We have $AB\ne0$, $\gcd(A,B)=1$ and $p\equiv1\pmod8$. In
$\mathbf F_p$ put $i=A/B$. Then $i^2=-1$. Quadratic reciprocity
gives $(B/p)=1$: the 2-part and sign are squares, and for every odd
prime $q\mid B$ one has $(p/q)=(A^2/q)=1$ and $p\equiv1\pmod4$.
Likewise, using Jacobi symbols if $A+B$ is composite,
\[
 \left(\frac{A+B}{p}\right)
 =\left(\frac{p}{|A+B|}\right)
 =\left(\frac{2}{|A+B|}\right)=1.
\]
Here $p\equiv2A^2\pmod {|A+B|}$, and $A+B\equiv1\pmod8$;
absolute value does not change the symbol of 2. Thus $2+2i$ is a
square. Its product with $2-2i$ is 8, a square modulo $p$, so the
latter is a square as well. Choose their square roots $u,v$.
Then $uv=\pm2\sqrt2$, and
$(u\pm v)^2/4=1\pm\sqrt2$. Both $1+\sqrt2$ and $1-\sqrt2$
are therefore squares, for either choice of the root of 2.

The three differences of the abscissa of $R_0$ from the roots
$0,d,-d$ are
\[
 1+\rho,\qquad 2(\rho-1)=2/(1+\rho),\qquad4.
\]
They are nonzero squares modulo $p$. The elementary full-two-torsion
descent criterion for $y^2=(x-e_1)(x-e_2)(x-e_3)$ therefore puts the
point in $2E(\mathbf F_p)$. This criterion follows directly from the
halving formulas: choose square roots of the three differences with
product $-y$ and use their pairwise products to construct the abscissa
of a half-point. The same calculation applies to the conjugate.

Finally a finite elliptic group with full rational two-torsion is
$\Z/a\Z\times\Z/b\Z$ with $a\mid b$ and both $a,b$ even.
The exponent of its doubled subgroup is $b/2$, which divides
$ab/4=\#E(\mathbf F_p)/4$.
\end{proof}

The inherited simultaneous-prime construction already imposes the
residue conditions in this lemma. Thus the existing conjugate order
moduli may be used for $R_0$, not merely for $P=2R_0$. No additional
prime-distribution theorem is required. The exact checker supplies
33 prime examples and a wrong-residue counterexample; these finite
calculations supplement, rather than replace, the proof above.

\begin{thebibliography}{9}
\bibitem{VB} J. von Buhren, \emph{Borne de hauteur semi-effective pour
le probl\`eme de Mordell--Lang dans une vari\'et\'e ab\'elienne},
Journal de Th\'eorie des Nombres de Bordeaux \textbf{29} (2017),
307--320, Theorem~1.1, printed p.~308.
\url{https://www.numdam.org/article/JTNB_2017__29_1_307_0.pdf}.
\bibitem{Trace} \emph{Trace heights on quadratic twists with $j=1728$},
project companion, 13 November 2026,
\nolinkurl{trace_on_all_points_20261113_v25.tex}.
\bibitem{Quartic} \emph{A quartic in 204 variables from recovered elliptic differences},
project companion, 14 November 2026,
\nolinkurl{quartic204_20261114_v27.tex}.
\end{thebibliography}
\end{document}
